% GENERATED FILE. Edit canonical lessons, then run npm run generate:books.
% canonical-source-hash: fnv1a64:efb0ad1ffb1aaf0d
% canonical-lessons: SA-C383-sankalpah, SA-C383-akanksa, SA-C383-akanksati, SA-C383-abhilasah, SA-C383-uttarottaram

\chapter{Setting a Goal}
\label{ch:sa-setting-a-goal}
\hlchaptermodality{383}

\begin{chapteropening}
\textbf{By the end of this chapter:} \emph{I can say what I mean to achieve and what it will take.}

Complete the last lesson of chapter 383: I can say what I mean to achieve and what it will take.
\end{chapteropening}

\section[\texorpdfstring{\sk{संकल्पः} (saṅkalpaḥ)}{संकल्पः (saṅkalpaḥ)}]{\sk{संकल्पः} (saṅkalpaḥ) — a resolve}

\label{lesson:SA-C383-sankalpah}



\begin{quote}
Before the new one: say the Sanskrit for rests, then the Sanskrit for receives with honour.
\end{quote}

\subsection*{You'll want to know: \sk{संकल्पः}}
\textbf{\sk{संकल्पः}} — \emph{saṅkalpaḥ} — \textquotedblleft{}a resolve\textquotedblright{}.

\begin{quote}
\sk{यदि} \sk{संकल्पः} \sk{दृढः}, \sk{तर्हि} \sk{लक्ष्यं} \sk{प्राप्नोति।}
\end{quote}

\begin{quote}
\emph{yadi saṅkalpaḥ dṛḍhaḥ, tarhi lakṣyaṃ prāpnoti.}
\end{quote}

\begin{quote}
If the resolve is firm, one reaches the goal.
\end{quote}

A masculine noun.

\begin{cousinweb}
\textbf{\sk{संकल्पः}} (\emph{saṅkalpaḥ}) is built on \emph{sam} and the root √kḷp, \textquotedblleft{}to be fit, to form\textquotedblright{}; MW gives \emph{saṃ-kalpa} \textquotedblleft{}will, volition, purpose, a definite intention\textquotedblright{} (MW).
\end{cousinweb}

\begin{grammarlens}[title={What you've built}]
One more word for saying what happens if.
\end{grammarlens}

\subsection*{Your turn}
\begin{itemize}
\raggedright
  \item \emph{Say it:} \emph{saṅkalpaḥ}
  \item \emph{Say it:} \emph{saṅkalpaḥ}, once more
  \item \emph{Say it:} \emph{yadi saṅkalpaḥ dṛḍhaḥ, tarhi lakṣyaṃ prāpnoti.} — If the resolve is firm, one reaches the goal.
  \item \emph{Recall:} say the Sanskrit for rests, then the Sanskrit for receives with honour, then say \emph{saṅkalpaḥ} again
\end{itemize}

\begin{tcolorbox}[breakable,colback=teal!4,colframe=teal!35!black,title={Before you move on}]
What does \sk{संकल्पः} mean? (\textquotedblleft{}A resolve\textquotedblright{}.) Say it once more.
\end{tcolorbox}
\section[\texorpdfstring{\sk{आकांक्षा} (ākāṅkṣā)}{आकांक्षा (ākāṅkṣā)}]{\sk{आकांक्षा} (ākāṅkṣā) — an ambition; a longing}

\label{lesson:SA-C383-akanksa}



\begin{quote}
Before the new one: say the Sanskrit for receives with honour, then the Sanskrit for a resolve.
\end{quote}

\subsection*{You'll want to know: \sk{आकांक्षा}}
\textbf{\sk{आकांक्षा}} — \emph{ākāṅkṣā} — \textquotedblleft{}an ambition; a longing\textquotedblright{}.

\begin{quote}
\sk{यदि} \sk{आकांक्षा} \sk{महती}, \sk{तर्हि} \sk{प्रयत्नः} \sk{अपि} \sk{महान्} \sk{भवेत्।}
\end{quote}

\begin{quote}
\emph{yadi ākāṅkṣā mahatī, tarhi prayatnaḥ api mahān bhavet.}
\end{quote}

\begin{quote}
If the ambition is great, the effort should be great too.
\end{quote}

A feminine noun.

\begin{cousinweb}
\textbf{\sk{आकांक्षा}} (\emph{ākāṅkṣā}) is built on \emph{ā} and the root √kāṅkṣ, \textquotedblleft{}to desire\textquotedblright{}; MW gives \emph{ā-kāṅkṣā} \textquotedblleft{}desire, wish\textquotedblright{} (MW).
\end{cousinweb}

\begin{grammarlens}[title={What you've built}]
One more word for saying what happens if.
\end{grammarlens}

\subsection*{Your turn}
\begin{itemize}
\raggedright
  \item \emph{Say it:} \emph{ākāṅkṣā}
  \item \emph{Say it:} \emph{ākāṅkṣā}, once more
  \item \emph{Say it:} \emph{yadi ākāṅkṣā mahatī, tarhi prayatnaḥ api mahān bhavet.} — If the ambition is great, the effort should be great too.
  \item \emph{Recall:} say the Sanskrit for receives with honour, then the Sanskrit for a resolve, then say \emph{ākāṅkṣā} again
\end{itemize}

\begin{tcolorbox}[breakable,colback=teal!4,colframe=teal!35!black,title={Before you move on}]
What does \sk{आकांक्षा} mean? (\textquotedblleft{}An ambition; a longing\textquotedblright{}.) Say it once more.
\end{tcolorbox}
\section[\texorpdfstring{\sk{आकांक्षति} (ākāṅkṣati)}{आकांक्षति (ākāṅkṣati)}]{\sk{आकांक्षति} (ākāṅkṣati) — longs for}

\label{lesson:SA-C383-akanksati}



\begin{quote}
Before the new one: say the Sanskrit for a resolve, then the Sanskrit for an ambition.
\end{quote}

\subsection*{You'll want to know: \sk{आकांक्षति}}
\textbf{\sk{आकांक्षति}} — \emph{ākāṅkṣati} — \textquotedblleft{}longs for\textquotedblright{}.

\begin{quote}
\sk{यदि} \sk{सः} \sk{ज्ञानम्} \sk{आकांक्षति}, \sk{तर्हि} \sk{प्रतिदिनं} \sk{पठेत्।}
\end{quote}

\begin{quote}
\emph{yadi saḥ jñānam ākāṅkṣati, tarhi pratidinaṃ paṭhet.}
\end{quote}

\begin{quote}
If he longs for knowledge, he should study every day.
\end{quote}

Past (he …): \sk{सः} \sk{ज्ञानम्} \sk{आकांक्षितवान्}, \emph{saḥ jñānam ākāṅkṣitavān}.

\begin{cousinweb}
\textbf{\sk{आकांक्षति}} (\emph{ākāṅkṣati}) is built on \emph{ā} and the root √kāṅkṣ, \textquotedblleft{}to desire, to long for\textquotedblright{}; MW gives \emph{ā-kāṅkṣā} \textquotedblleft{}desire, wish\textquotedblright{}, and KRM gives \emph{kāṅkṣita} for the root (MW).
\end{cousinweb}

\begin{grammarlens}[title={What you've built}]
One more word for saying what happens if.
\end{grammarlens}

\subsection*{Your turn}
\begin{itemize}
\raggedright
  \item \emph{Say it:} \emph{ākāṅkṣati}
  \item \emph{Say it:} \emph{ākāṅkṣati}, once more
  \item \emph{Say it:} \emph{yadi saḥ jñānam ākāṅkṣati, tarhi pratidinaṃ paṭhet.} — If he longs for knowledge, he should study every day.
  \item \emph{Recall:} say the Sanskrit for a resolve, then the Sanskrit for an ambition, then say \emph{ākāṅkṣati} again
\end{itemize}

\begin{tcolorbox}[breakable,colback=teal!4,colframe=teal!35!black,title={Before you move on}]
What does \sk{आकांक्षति} mean? (\textquotedblleft{}Longs for\textquotedblright{}.) Say it once more.
\end{tcolorbox}
\section[\texorpdfstring{\sk{अभिलाषः} (abhilāṣaḥ)}{अभिलाषः (abhilāṣaḥ)}]{\sk{अभिलाषः} (abhilāṣaḥ) — a heartfelt wish}

\label{lesson:SA-C383-abhilasah}



\begin{quote}
Before the new one: say the Sanskrit for an ambition, then the Sanskrit for longs for.
\end{quote}

\subsection*{You'll want to know: \sk{अभिलाषः}}
\textbf{\sk{अभिलाषः}} — \emph{abhilāṣaḥ} — \textquotedblleft{}a heartfelt wish\textquotedblright{}.

\begin{quote}
\sk{यदि} \sk{तव} \sk{अभिलाषः} \sk{शुद्धः}, \sk{तर्हि} \sk{सः} \sk{पूर्णः} \sk{भविष्यति।}
\end{quote}

\begin{quote}
\emph{yadi tava abhilāṣaḥ śuddhaḥ, tarhi saḥ pūrṇaḥ bhaviṣyati.}
\end{quote}

\begin{quote}
If your wish is pure, it will be fulfilled.
\end{quote}

A masculine noun.

\begin{cousinweb}
\textbf{\sk{अभिलाषः}} (\emph{abhilāṣaḥ}) is built on \emph{abhi}, \textquotedblleft{}towards\textquotedblright{}, and the root √laṣ, \textquotedblleft{}to desire\textquotedblright{}; MW gives \emph{abhi-lāṣa} \textquotedblleft{}desire, wish, affection\textquotedblright{} (MW).
\end{cousinweb}

\begin{grammarlens}[title={What you've built}]
One more word for saying what happens if.
\end{grammarlens}

\subsection*{Your turn}
\begin{itemize}
\raggedright
  \item \emph{Say it:} \emph{abhilāṣaḥ}
  \item \emph{Say it:} \emph{abhilāṣaḥ}, once more
  \item \emph{Say it:} \emph{yadi tava abhilāṣaḥ śuddhaḥ, tarhi saḥ pūrṇaḥ bhaviṣyati.} — If your wish is pure, it will be fulfilled.
  \item \emph{Recall:} say the Sanskrit for an ambition, then the Sanskrit for longs for, then say \emph{abhilāṣaḥ} again
\end{itemize}

\begin{tcolorbox}[breakable,colback=teal!4,colframe=teal!35!black,title={Before you move on}]
What does \sk{अभिलाषः} mean? (\textquotedblleft{}A heartfelt wish\textquotedblright{}.) Say it once more.
\end{tcolorbox}
\section[\texorpdfstring{\sk{उत्तरोत्तरम्} (uttarottaram)}{उत्तरोत्तरम् (uttarottaram)}]{\sk{उत्तरोत्तरम्} (uttarottaram) — more and more; step after step}

\label{lesson:SA-C383-uttarottaram}



\begin{quote}
Before the new one: say the Sanskrit for longs for, then the Sanskrit for a heartfelt wish.
\end{quote}

\subsection*{You'll want to know: \sk{उत्तरोत्तरम्}}
\textbf{\sk{उत्तरोत्तरम्}} — \emph{uttarottaram} — \textquotedblleft{}more and more; step after step\textquotedblright{}.

\begin{quote}
\sk{यदि} \sk{भवान्} \sk{यतते}, \sk{तर्हि} \sk{उत्तरोत्तरं} \sk{कुशलः} \sk{भविष्यति।}
\end{quote}

\begin{quote}
\emph{yadi bhavān yatate, tarhi uttarottaraṃ kuśalaḥ bhaviṣyati.}
\end{quote}

\begin{quote}
If you keep trying, you will grow more and more skilled.
\end{quote}

\begin{cousinweb}
\textbf{\sk{उत्तरोत्तरम्}} (\emph{uttarottaram}) is built from \emph{uttara}, \textquotedblleft{}higher, later\textquotedblright{}, said twice; MW gives \emph{uttarottaram} \textquotedblleft{}higher and higher, more and more, in constant continuation\textquotedblright{} (MW).
\end{cousinweb}

\begin{grammarlens}[title={What you've built}]
One more word for saying what happens if.
\end{grammarlens}

\subsection*{Your turn}
\begin{itemize}
\raggedright
  \item \emph{Say it:} \emph{uttarottaram}
  \item \emph{Say it:} \emph{uttarottaram}, once more
  \item \emph{Say it:} \emph{yadi bhavān yatate, tarhi uttarottaraṃ kuśalaḥ bhaviṣyati.} — If you keep trying, you will grow more and more skilled.
  \item \emph{Recall:} say the Sanskrit for longs for, then the Sanskrit for a heartfelt wish, then say \emph{uttarottaram} again
\end{itemize}

\begin{tcolorbox}[breakable,colback=teal!4,colframe=teal!35!black,title={Before you move on}]
What does \sk{उत्तरोत्तरम्} mean? (\textquotedblleft{}More and more; step after step\textquotedblright{}.) Say it once more.
\end{tcolorbox}
